\section{Day 2: Algebra Preliminaries; Algebraic Numbers and Integers (Sep.\ 10, 2026)}
We set some conventions for this class. Unless specified otherwise, a ring will always refer to a commutative ring with unity.
\begin{definition}[Preliminaries] \label{def:class_conventions}
    Let $R$ be a ring.
    \begin{enumerate}[(i)]
        \item $R$ is an \textit{integral domain} if it admits no zero divisors, i.e., whenever $ab = 0$ for $a, b \in R$, then $a = 0$ or $b = 0$.
        \item We will write $a \mid b$ if ``$a$ divides $b$'' for $a, b \in R$ if there exists some $c \in R$ such that $b = ac$.
        \item An integral domain $R$ is a \textit{principal ideal domain} (PID) if every ideal is principal, i.e., generated by a single element. Specifically, if $I \subset R$ is an ideal, then there exists $a \in R$ such that $I = \left<a\right>$.
        \item An element $u \in R$ is a \textit{unit} if it has a multiplicative inverse. An element $p \in R$ is \textit{irreducible} if it is not a unit, and whenever $p = ab$, either $a$ or $b$ is a unit. An element $p \in R$ is \textit{prime} if, whenever $p \mid ab$, then $p \mid a$ or $p \mid b$.
        \item An integral domain $R$ is a \textit{unique factorization domain} (UFD) if every non-unit $x \in R$ is uniquely (up to permutation and multiplication by units) a product of irreducible elements.
    \end{enumerate}
\end{definition}

\begin{proposition} \label{prop:UFD_irreducible_prime}
    If $R$ is a UFD, then irreducibles are primes, and primes are irreducible.
\end{proposition}
Note that primes are always irreducible. You can reference my MAT347 notes for a proof on this.

\begin{proposition} \label{prop:quotient_by_prime_ideal_iff_integral_domain} % man. what the fuck do i name this. xd
    Let $R$ be a ring. Then $p \in R$ is prime if and only if $R/\left<p\right>$ is an integral domain.
\end{proposition}
\begin{proof}
    Pick $\ol a, \ol b \in R/\left<p\right>$ such that $\ol{ab} = 0$ in $R/\left<p\right>$. Then $ab \in \left<p\right>$, so $p \mid ab$, where $p \mid a$ or $p \mid b$ as $p$ is prime, so $\ol a = 0$ or $\ol b = 0$. The converse direction is obtained by reversing this proof.
\end{proof}

We now continue writing out the definitions.
\begin{definition}[Preliminaries, cont.] \label{def:class_conventions_2}
    Let $R$ be a ring.
    \begin{enumerate}[(i)] \setcounter{enumi}{5}
        \item An integral domain $R$ is a \textit{Euclidean domain} (ED) if there exists a map $v \colon R \setminus \{0\} \to \ZZ_{\geq 0}$ such that, for all $a, b \in R$ with $b \neq 0$, there exists $q, r \in R$ such that $a = bq + r$ with $r = 0$ or $v(r) < v(b)$. In addition, for all nonzero $a, b \in R$, $v$ must satisfy $v(a) \leq v(ab)$.% {\color{airforceblue} The two properties may be thought of as ``division with remainder''}
    \end{enumerate}
\end{definition}

\begin{theorem} \label{thm:ed_is_pid_is_ufd}
    Let $R$ be an integral domain. If $R$ is a Euclidean domain, then it is a PID. If $R$ is a PID, then it is a UFD.
\end{theorem}
Once again, the proof is in my MAT347 notes. Let us give some examples of this.
\begin{itemize}
    \item $R = \ZZ$ is a Euclidean domain with $v(n) = \abs{n}$.
    \item If $F$ is a field, then $F[x]$ is a Euclidean domain with $v(f(x)) = \deg f$.
    \item Let $n$ be an integer, and let $R = \ZZ/n$ be the integers mod $n$. $R$ is an integral domain if and only if $n$ is prime; in fact, if $n$ is prime, then $\ZZ/n$ is a field. For a quick proof of this fact, let $x \in \ZZ$; then $\ol x \in \ZZ/n$ is a unit if and only if $x$ is coprime to $n$. If $p$ is rpime, then $(\ZZ/p)^\times$ is a group of order $p-1$.
\end{itemize}

\begin{theorem} \label{thm:prime_number_properties}
    Let $p \in \NN$ be a prime number.
    \begin{enumerate}[(i)]
        \item (Fermat's little theorem) If $a \in \ZZ$ is coprime to $p$, then $a^{p-1} \equiv 1 \pmod{p}$.
        \item $(\ZZ/p)^\times$ is a cyclic group.
    \end{enumerate}
\end{theorem}

\newpage
\begin{corollary} \label{cor:primitive_root}
    Let $p \in \NN$ be a prime number. Then there exists $g \in \ZZ$, called the \textit{primitive root modulo $p$}, such that \begin{parlist}
        \item if $a$ is coprime to $p$, then $a \equiv g^r \!\! \pmod{p}$ for some $r \geq 0$, and
        \item if $r \in \ZZ$ such that $g^r \equiv 1 \!\! \pmod{p}$, then $p-1 \mid r$.
    \end{parlist}
\end{corollary}

\begin{definition} \label{def:algebraic_number_and_integer}
    An \textit{algebraic number} $\alpha$ is a complex number $\alpha \in \CC$ for which there exists $a_0, \dots, a_{n-1} \in \QQ$ (with $n \geq 1$) such that
    \[ \alpha^n + a_{n-1} \alpha^{n-1} + \dots + a_1 \alpha + a_0 = 0, \]
    i.e., $p(\alpha) = 0$ for some monic $p(x) \in \QQ[x]$. An \textit{algebraic integer} is a complex number $\alpha \in \CC$ such that there exists $a_0, \dots, a_{n-1} \in \ZZ$ (with $n \geq 1$) such that
    \[ \alpha^n + a_{n-1} \alpha^{n-1} + \dots + a_1 \alpha + a_0 = 0, \]
    i.e., $p(\alpha) = 0$ for some monic $p(x) \in \ZZ[x]$. Algebraic numbers are denoted by $\QQ^{\opname{alg}}$ or $\ol\QQ$, and algebraic integers are denoted by $\ZZ^{\opname{alg}}$ or $\ol\ZZ$. The bar denotes algebraic closure.
\end{definition}
In particular, algebraic numbers form a field, and algebraic integers form a ring. This is not obvious from the definition, but we will take it as true for now. Let us give some examples.
\begin{itemize}
    \item Any $\alpha \in \QQ$ is an algebraic number, since we may pick $p(x) = x - \alpha = 0$. Similarly, any $n \in \ZZ$ is an algebraic integer per the exact same construction.
    \item $\ZZ^{\opname{alg}} \subset \QQ^{\opname{alg}}$.
    \item $\sqrt{-5} \in \ZZ^{\opname{alg}}$, since $(\sqrt{-5})^2 + 5 = 0$.
    \item $\omega = \frac{1}{2} (-1 + \sqrt{-3})$ an algebraic integer, since $\omega^2 + \omega + 1 = 0$.
    \item For $n \geq 1$, we have that $\zeta = \exp(2\pi i /n) \in \ZZ^{\opname{alg}}$, since $\zeta^n - 1 = 0$. $\zeta$ is the \textit{primitive $n^{\mathrm{th}}$ root of unity}.
    \item Let $\alpha \in \QQ^{\opname{alg}}$ and $r \in \QQ$. Then $r \alpha \in \QQ^{\opname{alg}}$. We will see later that, for all $\alpha \in \QQ^{\opname{alg}}$, there exists some $r \in \ZZ$ such that $r \alpha \in \ZZ^{\opname{alg}}$.
    \item $\QQ^{\opname{alg}}$ is countable, since it is the union of the (finitely many) roots of polynomials in $\QQ[x]$, of which are countable. $\CC$ is not countable, however, so $\QQ^{\opname{alg}} \subsetneq \CC$.
    \item $e, \pi \notin \QQ^{\opname{alg}}$.
\end{itemize}

Another way to characterize $\QQ^{\opname{alg}}$ and $\ZZ^{\opname{alg}}$ is to consider subsets $V \subset \CC$ such that $V$ is a vector space (respectively, a finitely generated abelian group) over $\QQ$. As an example, $V = \QQ$ is trivially a vector space over $\QQ$, and $V = \{r + s \sqrt{-5} \mid r, s \in \QQ\}$ is one as well. In particular, it contains the finitely generated abelian group $\{r + s \sqrt{-5} \mid r, s \in \ZZ\}$.
\begin{lemma} \label{lem:check_for_algebraic}
    Let $\alpha \in \CC$. Then
    \begin{enumerate}[(i)]
        \item $\alpha \in \QQ^{\opname{alg}}$ if and only if there exists a nonzero finite dimensional vector space $V \subset \CC$ over $\QQ$ such that $\alpha V \subset V$.
        \item $\alpha \in \ZZ^{\opname{alg}}$ if and only if there exists a finitely generated abelian group $M \subset \CC$ such that $\alpha M \subset M$.
    \end{enumerate}
\end{lemma}
\begin{proof}
    We start by proving (i). In the forwards direction, suppose $\alpha \in \QQ^{\opname{alg}}$; then $\alpha^n + a_{n-1} \alpha^{n-1} + \dots + a_0 = 0$. Let
    \[ V = \{r_0 + r_1\alpha + \dots + r_{n-1}\alpha^{n-1} \mid r_0, \dots, r_{n-1} \in \QQ\}, \]
    which we readily see as finite-dimensional as a vector space over $\QQ$. It follows that $\alpha V \subset V$, since
    \begin{align*}
        \alpha(r_0 + r_1 \alpha + \dots + r_{n-1}\alpha^{n-1}) &= r_0 \alpha + r_1 \alpha^2 + \dots + r_{n-1} \alpha^n \\
        &= r_0\alpha + r_1\alpha^2 + \dots + r_{n-1}(-a_{n-1}\alpha^{n-1} - \dots - a_0) \in V.
    \end{align*}
    Conversely, suppose $V \subset \CC$ is finite-dimensional with $\alpha V \subset V$. Let $v_1, \dots, v_d$ be a basis of $V$ over $\QQ$. For each $i$, we see that
    \[ \alpha v_i = \sum_{j=1}^d r_{ij} v_j, \qquad r_{ij} \in \QQ; \]
    i.e., $v \mapsto \alpha v$ is represented by a matrix $R \coloneq (r_{ij})$ in this basis. Specifically, $R$ is a $d \times d$ matrix with entries in $\QQ$, which we may think of as a complex matrix:
    \[ R \begin{pmatrix} v_1 \\ \vdots \\ v_d \end{pmatrix} = \begin{pmatrix} \sum_{j=1}^d r_{ij} v_j \\ \vdots \\ \sum_{j=1}^d r_{dj} v_j \end{pmatrix} = \alpha \begin{pmatrix} v_1 \\ \vdots \\ v_d \end{pmatrix}. \]
    It follows that $\alpha \in \CC$ is an eigenvalue of $R$, whence $\alpha$ is a root of the characteristic polynomial $\det (R - I x) \in \QQ[x]$, so it has to be an algebraic number.

    The proof for (ii), i.e., the algebraic integers case, follows analogously. Note that we need to know that the finitely generated abelian group is free; per the classification theorem (of finitely generated abelian groups), every such $V$ can be written in terms of its free part and torsion; but there is no torsion, so we may reiterate our earlier argument as desired.
\end{proof}

\begin{proposition} \label{prop:algebraics_form_subfield_subring}
    $\QQ^{\opname{alg}} \subset \CC$ is a subfield, and $\ZZ^{\opname{alg}} \subset \CC$ is a subring.
\end{proposition}
\begin{proof}
    Let $\alpha, \beta \in \QQ^{\opname{alg}}$. Then, per \ref{lem:check_for_algebraic}, there exist nonzero finite-dimensional subspaces $V, W \subset \CC$ such that $\alpha V \subset V$ and $\beta W \subset W$. Consider
    \[ V \cdot W = \{v_1w_1 + \dots v_tw_t \mid v_j \in V, w_j \in W\}; \]
    observing that $v_1, \dots, v_n$ span $V$ and $w_1, \dots, w_m$ span $W$, we simply have that $v_i w_j$ for $1 \leq i \leq n$ and $1 \leq j \leq m$ span $V \cdot W$. Thus, $V \cdot W$ is finite-dimensional and nonempty.

    Now, for any $v \in V$ and $w \in W$, we have that $\alpha \cdot (vw) = (\alpha v) \cdot w \in V \cdot W$, and $\beta \cdot (vw) = v \cdot (\beta w) \in V \cdot W$. By linearity for $z \in V \cdot W$, we have $\alpha z \in V \cdot W$ and $\beta z \in V \cdot W$, whence $(\alpha + \beta) (V \cdot W) \subset V \cdot W$ and $(\alpha \cdot \beta) (V \cdot W) \subset V \cdot W$, so $\alpha + \beta$ and $\alpha \beta$ are both algebraic numbers per \ref{lem:check_for_algebraic}. An analogous argument holds for $\alpha, \beta \in \ZZ^{\opname{alg}}$.

    It remains to show that nonzero $\alpha \in \QQ^{\opname{alg}}$ have inverses, i.e., $\alpha\inv \in \QQ^{\opname{alg}}$. Indeed, the ``multiplying by $\alpha$'' map, $\alpha \colon V \to V$, is injective, so it is isomorphic, and so there exists $\alpha\inv$ as the inverse map to $\alpha$. In particular, it preserves $V$. Alternatively, we may write
    \[ \alpha^n + a_{n-1} \alpha^{n-1} + \dots + a_0 = 0 \implies \alpha \cdot \left(-\frac{1}{\alpha_0}\right) (\alpha^{n-1} + a_{n-1} \alpha^{n-2} + \dots + a_1) = 1, \]
    which is a more explicit construction for $\alpha\inv$.
\end{proof}

\newpage
We now look at more examples.
\begin{itemize}
    \item Every element of $\ZZ[\sqrt{-5}]$ is an algebraic integer. Indeed, every $\alpha \in \ZZ[\sqrt{-5}]$ is of the form $a + b \sqrt{-5}$ with $a, b \in \ZZ$, and there is clearly an easy monic quadratic $p(x) \in \ZZ[x]$ such that $p(\alpha) = 0$.
    \item $\alpha = \sqrt{2} + \sqrt{3} \in \ZZ^{\opname{alg}}$, since it satisfies $\alpha^4 - 10 \alpha^2 + 1 = 0$.
\end{itemize}

\begin{proposition} \label{prop:minimal_poly}
    Let $\alpha \in \QQ^{\opname{alg}}$. Then there exists a unique monic irreducible polynomial $f(x) \in \QQ[x]$ such that $f(\alpha) = 0$. Moreover, if $g(x) \in \QQ[x]$ also satisfies $g(\alpha) = 0$, then $f \mid g$.
\end{proposition}
\begin{proof}
    Consider $I = \{h(x) \in \QQ[x] \mid h(\alpha) = 0\} \neq \emptyset$ (nonempty, as $\alpha$ is algebraic). $I$ is clearly an ideal, %since the sum of two polynomials with $\alpha$ as a root also has $\alpha$ as a root.
    and $\QQ[x]$ is a PID, so we may write $I = \left<f\right>$ for some $f$. Since $f$ is uniquely determined up to units, we may divide any choice of $f$ by its leading coefficient to let $f$ be monic. It remains to check that $f$ is irreducible. Let $f = gh$; then $0 = f(\alpha) = g(\alpha) h(\alpha)$. It must be that one of $g(\alpha)$ or $h(\alpha)$ is zero, with the other being in $I$. Without loss of generality, let $g(x) \in I$; but then $f \mid g$, so $\deg g \geq \deg f$, whence $h$ is a unit.
\end{proof}